% Source-audited historical slice. Chronology controls the major order;
% functional type controls the order within each period.
\part{Objects Used in Earlier Roman Practice}

\begin{center}
\fbox{\begin{minipage}{0.91\linewidth}
\small
\textbf{Selective historical plate set for review.}
These six entries test the chronological, source-note, and pencil-plate
method. They are not a complete inventory of objects used before 1962, and
none is presented as an instruction for the 1962 rites.
\end{minipage}}
\end{center}

\section{Early Christian and Medieval Forms}
\DictionaryPlateHeader{Historical plate H1}{Reservation and Liturgical Service}
  {not to common scale}
\DictionaryGraphic{3.55in}
  {liturgy/roman-rite/1962/reference/roman-sanctuary-dictionary/shared/artwork/pencil/RSD-HIST-001-early-medieval-objects.png}

\DictionaryThreeCellRow
{\DictionaryCell{obj-hist-flabellum}{Liturgical fan}{flabellum}
  {}{A hand fan used near the altar to keep insects from the sacred species.}
  {HISTORICAL WESTERN USE}{Its practical care for the Eucharist also became a
  ceremonial sign of honor.}}
{\DictionaryCell{obj-hist-eucharistic-dove}{Eucharistic dove}{columba eucharistica}
  {}{A dove-shaped outer vessel or pyx suspended above the altar for
  Eucharistic reservation.}{HISTORICAL RESERVATION FORM}{The dove made the
  association with the Holy Spirit visible.}}
{\DictionaryCell{obj-hist-eucharistic-tower}{Eucharistic tower}{turris eucharistica}
  {}{A tower-shaped receptacle for Eucharistic reservation, placed in an
  armarium, near the altar, or in some examples suspended.}
  {HISTORICAL RESERVATION FORM}{}}

\section{Early Modern Continuities}
This first audited slice adds no separate object merely because its ornament
changed. Chasuble cut, vessel decoration, and architectural style remain
variant questions until a source establishes a materially different function
or construction.

\section{Twentieth-Century Reform before 1962}
\DictionaryPlateHeader{Historical plate H2}{Forms Absent from the Revised Holy Week}
  {not to common scale}
% The source is 1024 px wide. Keep its largest review placement below 3.42 in
% so the trimmed historical study remains at or above 300 effective dpi.
\DictionaryGraphic[trim=0 300 0 80,clip]{3.35in}
  {liturgy/roman-rite/1962/reference/roman-sanctuary-dictionary/shared/artwork/pencil/RSD-HIST-002-pre-1955-objects.png}

\DictionaryThreeCellRow
{\DictionaryCell{obj-hist-folded-chasuble}{Folded chasuble}{planeta plicata}
  {}{A chasuble worn folded up by a sacred minister in place of dalmatic or
  tunicle on specified penitential occasions before reform.}
  {HISTORICAL IN THE REVISED HOLY-WEEK ROUTE}{}}
{\DictionaryCell{obj-hist-broad-stole}{Broad stole}{stola latior}
  {}{A broad diagonal band assumed by the deacon after laying aside the folded
  chasuble in the older ceremonial sequence.}
  {HISTORICAL IN THE REVISED HOLY-WEEK ROUTE}{}}
{\DictionaryCell{obj-hist-triple-candle}{Triple candle on reed}{arundo cum tribus candelis}
  {}{A reed-like staff bearing three candles, successively lighted during the
  older Easter Vigil entrance.}
  {SUPERSEDED BY THE REFORMED EASTER-VIGIL ENTRANCE}{}}

\section{Historical Source Notes}
\DictionaryEndnote{obj-hist-flabellum}{Joseph Braun, ``Flabellum,'' in
  \textit{The Catholic Encyclopedia}, vol.~6 (1909), records its practical
  purpose, materials, early witness, and continuation in the Latin Church to
  about the fourteenth century.}
\DictionaryEndnote{obj-hist-eucharistic-dove}{Joseph Braun, ``Dove,'' in
  \textit{The Catholic Encyclopedia}, vol.~5 (1909), describes the suspended
  dove as a Eucharistic vessel and distinguishes an immediate container from
  an outer vessel enclosing a pyx.}
\DictionaryEndnote{obj-hist-eucharistic-tower}{Joseph Braun, ``History of the
  Christian Altar,'' in \textit{The Catholic Encyclopedia}, vol.~1 (1907),
  describes dove and tower reservation forms and a thirteenth-century
  suspended tower at Arras.}
\DictionaryEndnote{obj-hist-folded-chasuble}{Herbert Thurston, ``Chasuble,'' in
  \textit{The Catholic Encyclopedia}, vol.~3 (1908), identifies the
  \textit{planetae plicatae}; Frederick R. McManus,
  \textit{The Rites of Holy Week} (1956), explicitly notes that folded
  chasubles are not worn in the revised rites.}
\DictionaryEndnote{obj-hist-broad-stole}{Thurston, ``Chasuble,'' describes the
  older folded-chasuble arrangement; McManus, \textit{Rites of Holy Week},
  records the revised vesture and thereby bounds this figure to the replaced
  route rather than the 1962 Holy Week.}
\DictionaryEndnote{obj-hist-triple-candle}{The pre-reform rubric is preserved in
  the contemporary \textit{Catholic Encyclopedia} discussion of the
  liturgical use of fire: ``\textit{Praeparetur arundo cum tribus candelis in
  summitate positis}.'' McManus, \textit{Rites of Holy Week}, describes the
  reformed entrance with the Paschal candle.}

\begin{center}
\footnotesize
\textbf{Boundary.} These six entries are a selective, source-audited slice,
not a claim that the historical inventory is complete. They are historical
comparanda and must not be read as directions for the 1962 rites.
\end{center}
